\scanpage{198}

\deleted{\unsure{donc}} (donc c'est la \inserted{$S$-}algèbre engendrée par les $U_{r,i}=U_r^{(i)}\quad 0\leq r\leq n-1,\ i\geq1$, soumis aux relations $U_{r,i}U_{r,j}=C_{ij}U_{r,i+j}$). \unsure{De $=$|Et} on a $\mathcal A_n\simeq S\{U_0,-,U_{n-1}\}[dU_0,-,dU_{n-1}]$

C'est immédiat, car $X_0,-,X_{\unsure{r|n}-1}$ se réconstituent par $U_0,-,U_{n-1}$ par

\[
X_0=U_0,\ X_1=U_1-U_0,\ --,\ X_{n-1}=U_{n-1}-U_{n-2}
\]

et on a $\mathcal A_n^0\simeq S\{X_0,-,X_{n-1}\}$. \unsure{Or}

Ceci montre l'unicité de $\varphi_n$, et qu'on devra avoir

\[
\begin{aligned}
\varphi_n(\underbrace{U_{r,i}^{\deleted{\unsure{(i)}}}}_{U_r^{(i)}})
 &=\varphi_n\bigl(\underbrace{\sigma_r^*}_{\substack{\Vert\,\text{déf}\\\mathcal A_*(\sigma_r)}}(X_0^{(i)})\bigr)
 =\sigma_r^*\bigl(\varphi(X_0^{(i)})\bigr).\\
 &=\sigma_r^*(\xi_i)\ :
\end{aligned}
\]

i.e.

\[
\varphi_n(U_{r\deleted{\unsure{,i}}}^{(i)})=\sigma_r^*(\xi_i)
=\biggl(\deleted{\unsure{A_r U}}\sum_j a_{ij}U_r^{(j)}+\Bigl(\sum b_{ij}U_r^{(j)}\Bigr)dU_r\biggr)
\]

Il y a \inserted{donc} un hom. d'algèbres \inserted{unique} $\mathcal A_n^0\overset{\varphi_n^0}{\unsure{-|\to}}\mathcal A_n$ tel que

\[
\varphi_n^0(U_r^{(i)})=\sigma_r^*(\xi_i)\overset{\text{déf}}{=}\xi_{r,i}
\]

car cela revient à vérifier les relations

\[
\deleted{\sigma_r^*(\xi_i)\sigma_r^*(}\qquad
\xi_{r,i}\xi_{r,j}=C_{ij}\xi_{r,i+j}
\]

i.e.

\[
\begin{aligned}
\underset{\substack{\Vert\\\sigma_r^*(\xi_i\xi_j)}}{\sigma_r^*(\xi_i)\sigma_r^*(\xi_j)}&=C_{ij}\sigma_r^*(\xi_{i+j})\\
\sigma_r^*(C_{ij}\xi_{i+j})&=C_{ij}\sigma_r^*(\xi_{i+j})\qquad\mathrm{OK}
\end{aligned}
\]

% archive sheet number: 197
